% 由 PDF 自动转换生成 v3

\chapter*{内容提要}
本书是为工程力学专业本科专业基础课程"流体力学"编写的教材，也可作为大学工科相关专 业和研究生学习流体力学的教材和参考书.内容包括流体力学的基本概念、原理、应用和一部分重 要的近代流体力学知识.全书共分 10 章.第 1-4 章是流体力学的基本原理，包括流体的物理性 质、流体运动学、流体动力学的基本原理、理想流体动力学 s 第 5 、 6 章是理想不可压缩流体动力学 的主要应用.包括理想不可压缩流体的二维无旋和有旋流动、水波动力学 s 第 7 章着重介绍气体动 力学基础 s 第 8-10 章是粘性流体动力学.包括粘性流体力学基础、湍流和边界层理论基础.书中 每章列举丰富的例题和提供大量的习题.书后给出部分习题的答案.认真学完本书后，读者将具备 进一步学习流体力学专门知识或着手研究和解决工程及自然界流动问题的扎实基础. 本书是 2006 年版《流体力学以第 2 版)的修订版，内容有更新和修改. 版权所有，侵权必究。侵权举报电话: 010-62782989 13701121933 图书在版编目 (CIP) 数据 流体力学/张兆 )1颐，崔桂香编著. --3 版.一北京 s 清华大学出版社， 2015 (高等院校力学教材)

\[ISBN 978-7-302-39785-4\]

1.${\rm 1}$流… II. ${\rm 1}$张… ${\rm 2}$崔… 阻.${\rm 1}$流体力学一高等学校一教材 IV. ( 035 中国版本图书馆 CIP 数据核字 (2015) 第 077173 号 责任编辑:石磊赵从棉 封面设计 2 傅瑞学 责任校对 z 赵丽敏 责任印制=沈露 出版发行=清华大学出版社 网址: http://www. tup. com. cn , http://www.wqbook.com 地 址 2 北京清华大学学研大厦 A 座 邮 编: 100084 社总机: 010-62770175 邮 购: 010-62786544 技稿与读者服务: 010-62776969 , c-service@tup. tsinghua. edu. cn 质量反馈: 010-62772015 , zhiliang@tup. tsinghua. edu. cn 印刷者 z 三洞市君旺印务有限公司 装订者:三饲市新茂装订有限公司 经 销 2 全国新华书店 开本: 170mmX240mm 印张: 28 版 次: 1999 年 2 月第 1 版 2015 年 7 月第 3 版 印数: 1-2000 定价: 49.00 元 产品编号: 051285-01 字 数: 556 千字 印 次: 2015 年 7 月第 I 次印刷

\chapter*{第3版前言}
本书向问世以来，被很多高等院校师生所采用，修订后的第 2 版又印刷 6 次，发 行1. 4 万余册. 为了进一步提高可读性，增加启发性 z 满足读者和任课教师对流体力学概念与基 本原理"加深理解，拓展思维，引领创新"的需要，本书进行了以下修改和补充 z (1)每章习题后增加了少量"思考题"，以引导学生思考和讨论，加深理解，拓展 思维\$ (2) 修改了部分书中附图以改善视觉效果; (3) 第 2 章 2.6 节补充了 "F${\rm o}$pple 定理"，以便更好地理解绕流问题中旋涡生成 的运动学机制; (4) 订正了原书中书写和印刷错误. 再次感谢广大读者在使用本教材过程中提出的宝贵意见和建议. 作者 2015 年 4 月于清华园

\chapter*{第2版前言}
本书第 1 版深得读者的厚爱，已印刷 5 次，发行 1 万余册.为了更好地满足广大 读者的需要，我们对原书作了修订.根据作者和同事们使用本教材的经验，在保留原 书风格的基础上，我们对原书作了如下的删节和补充. (1)尽量删节理想流体动力学中已经较少应用的内容. (2) 增加粘性流体力学的内容，把原有的粘性流体力学基础扩展为两章.粘性流 体力学基础中增加润滑问题的近似解法:把边界层理论基础单设一章，并增加可压缩 边界层的分析方法和特性. (3) 原书祸及涡动力学一章的内容，略加删节后，合并到运动学和理想流体 动力学. (4) 增加部分习题答案. (5) 订正了第 1 版正文和附录中的书写与印刷错误. 作者对广大教师和学生在使用本教材过程中提出的宝贵意见深表感谢. 作者 2006 年 5 月于北京清华园

\chapter*{第1版前言}
本书是为工程力学专业大学本科专业基础课程"流体力学"编写的教材，也可作 为大学工科专业和研究生学习流体力学的参考书. 本书是在作者多年讲授流体力学课程的讲稿基础上编写而成的.根据作者的教 学经验，流体力学作为一门专业基础课或技术基础课，必须要求学生对基本概念和原 理理解得十分准确和透彻;另一方面，要求学生学会应用流体力学理论的基本方法. 本书力求能够满足上述要求，在本书前几章讲述基本原理时，对重要概念和定理都列 举示范例题，这些例题可以在课堂讲授，也可以要求学生自学.俗话说，熟能生巧，经 过反复练习以后，就能将理论运用自如.流体力学是 一 门应用很广的基础课，自学或 初学的读者往往忽视基本理论的学习和运用理论的熟练程度.急于想去解决复杂的 实际问题，这样常常事倍功半.所以，我们建议初学流体力学的学生或自学的读者循 序渐进，扎扎实实掌握原理，认认真真练习，才能较快入门. 随着近代科学技术突飞猛进的发展，流体力学学科本身也在不断发展.在处理新 旧知识的更新方面.我们保留了流体力学的基本概念和原理部分.并尽量用现代观念 和方法来叙述.近代流体力学需要解决的问题越来越复杂，不论是几何外形或流动中 包含的物理、化学过程都越来越复杂.但是，不论运动怎样复杂，它的基本规律是共同 的.例如，流体运动的分析，控制流动的质量、动量和能量的输运方程是 一 样的.对于 这一部分内容不仅保留，还应当反复练习.对于概念的叙述和表达，我们尽量采用已 经比较成熟的近代数学方法.例如，全书用向量和张量的表达式描述流体运动学和动 力学.在流体力学的解题方法方面，本书保留和发展了经典流体力学中对于流动模型 的简化和流动问题的准确数学提法，这典方法对于用数值计算方法解决日益增多的 复杂流动问题仍然是基本的和重要的，例如，边界层模型、理想流体运动模型 E 线性化

W 第 1 版前言 方法和近代摄动法.至于经典流体力学中一些已经不常用的方法则予以删除.本书是 流体力学的入门课程，不包括流体力学全部专门知识.学完本书以后.有些学生将进 一步修习流体力学的专门课程，有的读者将以本课程知识为基础，通过参阅专门的书 刊和文献，着手解决实际流动问题f.无论对于哪一种读者，本书将为他们铺平顺利过 渡的道路. 全书的编排如下.第 1\textasciitilde{}4 章是流体力学的基本原理.所有学习流体力学的学生 都必须准确和完整地掌握.第 5\textasciitilde{}8 章是流体力学原理在各方面的应用.这一部分可 以根据专业的性质，有选择地重点讲授.例如，以水动力学为主的专业.可以对气体力 学部分有重点地宣讲.以空气动力学为主的专业可以对水动力学部分有重点地讲授. 第 9 章湍流和第 10 章涡动力学是近代流体力学中发展较快的部分，凡是有学时的专 业，我们建议安排学习这两章的基本内容. 本书第 1 ， 2.3.6.9 ， 10 章由张兆顺撰写.第 4.5 ， 7.8 章和附录以及全部习题由崔 桂香撰写，全部内容都经两人反复讨论和修改后定稿.清华大学 t 程力学系流体力学 教研组的部分教师租学住曾对书稿提出过宝贵的意见，张晓航同学描绘本书部分插 图.作者对他们表示深切的感谢. 作者 1998 年 7 月于北京清华园
